\section{Divisor-geometric null profiles and an exact certificate}
\label{sec:geometric-certificates}

The minimax identity is abstract.  We now show that its null directions
occur on exact divisor grids, rather than only in a continuum model.

\subsection{Prescribed finite samples and jets}

\begin{theorem}[Geometric Mellin-null profile]
\label{thm:geometric-null-profile}
Let $t_1,\ldots,t_q\in\R$ and $s_1,\ldots,s_q\geq1$, and put
\(
 D=\sum_{\ell=1}^q s_\ell
\).
Fix a prime $p$, set $\Delta=2\log p$, and write
\begin{equation}\label{eq:null-polynomial}
 Q(z)=\prod_{\ell=1}^q
 (z-e^{it_\ell\Delta})^{s_\ell}
 =\sum_{j=0}^Dc_jz^j.
\end{equation}
On the factor pairs
\begin{equation}\label{eq:prime-power-factor-pairs}
 r=p^D,\qquad n_j=p^j,\qquad m_j=p^{D-j},\qquad0\leq j\leq D,
\end{equation}
the scale nodes are
\begin{equation}\label{eq:prime-power-scale-nodes}
 u_j=\log(n_j/m_j)=(-D+2j)\log p.
\end{equation}
The nonzero complex profile $c=(c_0,\ldots,c_D)$ satisfies
\begin{equation}\label{eq:null-jet-identities}
 \left.\frac{d^k}{dt^k}\sum_{j=0}^Dc_je^{itu_j}
 \right|_{t=t_\ell}=0,
 \qquad 1\leq\ell\leq q,\quad0\leq k<s_\ell.
\end{equation}
\end{theorem}

\begin{proof}
The exponential polynomial of the profile factors exactly as
\begin{align}
 F_c(t)
 &=\sum_{j=0}^Dc_je^{it(-D+2j)\log p}
 \notag\\
 &=e^{-itD\log p}Q(e^{it\Delta}).
 \label{eq:null-profile-factorization}
\end{align}
At $t=t_\ell$, the last factor has a zero of order $s_\ell$, proving
\eqref{eq:null-jet-identities}.
\end{proof}

If the sample multiset is invariant under $t\mapsto-t$, with equal
multiplicities, then $Q$ has real coefficients.  If it also contains
$t=0$, then
\(
 \sum_jc_j=Q(1)=0
\).
Writing $c=c^+-c^-$, the positive and negative parts have the same mass.
After normalization they give two probability profiles on the same divisor
nodes with identical prescribed samples and jets.  Their masses in a
window $V$ differ whenever
\(
 \sum_jc_jV(u_j)\ne0
\).

\begin{corollary}[Zero jets are finite differences]
\label{cor:zero-jets-finite-difference}
For jets of orders $0,\ldots,K$ at $t=0$, take
\begin{equation}\label{eq:finite-difference-coefficients}
 Q(z)=(z-1)^{K+1},
 \qquad
 c_j=(-1)^{K+1-j}\binom{K+1}{j}.
\end{equation}
Then all scale moments through order $K$ vanish, while the response to a
window is its $(K+1)$st finite difference on the divisor grid.
\end{corollary}

\begin{remark}[Arithmetic support versus arithmetic coefficients]
The nodes \eqref{eq:prime-power-factor-pairs} are genuine factor pairs of
one integer, but the profile $c$ is constructed adversarially.  It is not
the restriction of $\Lambda(n)-1$.  Hence
\cref{thm:geometric-null-profile} proves a limitation of finite-data
recovery over coefficient classes; it is not a counterexample for the
actual prime residual.
\end{remark}

\subsection{The row \texorpdfstring{$81+2=83$}{81+2=83}}

We give a fully exact certificate whose target is prime.  Take
\begin{equation}\label{eq:certificate-row}
 h=2,\qquad r=3^4=81,\qquad r+h=83.
\end{equation}
The five factor pairs and normalized scale nodes are
\begin{equation}\label{eq:certificate-nodes}
 \begin{array}{c|ccccc}
 j&0&1&2&3&4\\
 \hline
 (n_j,m_j)&(1,81)&(3,27)&(9,9)&(27,3)&(81,1)\\
 x_j=u_j/\log3&-4&-2&0&2&4
 \end{array}.
\end{equation}
Define two probability profiles
\begin{equation}\label{eq:certificate-profiles}
 \mu_+=\frac18(1,0,6,0,1),
 \qquad
 \mu_-=\frac18(0,4,0,4,0).
\end{equation}

\begin{proposition}[Exact four-jet ambiguity]
\label{prop:exact-four-jet-certificate}
The profiles in \eqref{eq:certificate-profiles} have the same moments of
orders $0,1,2,3$:
\begin{equation}\label{eq:certificate-common-moments}
 (M_0,M_1,M_2,M_3)=(1,0,4,0).
\end{equation}
Consequently their Mellin transforms have identical derivatives through
order three at $t=0$.  If $V\in C_c(\R)$ satisfies
\begin{equation}\label{eq:certificate-window}
 V(0)=1,\qquad \supp V\subset(-\log3,\log3),
\end{equation}
then
\begin{equation}\label{eq:certificate-window-gap}
 \int V\,d\mu_+=\frac34,
 \qquad
 \int V\,d\mu_-=0.
\end{equation}
On the prime-target row \eqref{eq:certificate-row}, the corresponding
formal target-weighted gap is $(3/4)\log83$.
\end{proposition}

\begin{proof}
Direct substitution at the five integer nodes in
\eqref{eq:certificate-nodes} gives
\eqref{eq:certificate-common-moments}.  Equality of derivatives follows
because the $k$th derivative multiplies the $k$th moment by
$(i\log3)^k$.  The support condition in
\eqref{eq:certificate-window} makes the window-value vector
$v=(0,0,1,0,0)$, so \eqref{eq:certificate-window-gap} follows from the
central masses in \eqref{eq:certificate-profiles}.  Finally
$\Lambda(83)=\log83$ is common to every factor pair of $r=81$.
\end{proof}

The certificate also has an exact minimax value.

\begin{proposition}[Sharp five-node minimax radius]
\label{prop:sharp-five-node-radius}
On the nodes $x=(-4,-2,0,2,4)$, let $v=(0,0,1,0,0)$ and let
$\cP_3$ be the restrictions of polynomials of degree at most three.  Then
\begin{equation}\label{eq:sharp-five-node-radius}
 \dist_{\ell^\infty}(v,\cP_3)=\frac38.
\end{equation}
An optimal polynomial is
\begin{equation}\label{eq:certificate-best-polynomial}
 q(x)=\frac58-\frac{x^2}{16}.
\end{equation}
\end{proposition}

\begin{proof}
The vector
\begin{equation}\label{eq:certificate-annihilator}
 c=(1,-4,6,-4,1)
\end{equation}
annihilates $1,x,x^2,x^3$ on the five nodes.  Since
$\|c\|_1=16$ and $\langle c,v\rangle=6$, duality gives the lower bound
$6/16=3/8$.  At the five nodes the residual $v-q$ for
\eqref{eq:certificate-best-polynomial} alternates between $3/8$ and
$-3/8$, proving the matching upper bound.
\end{proof}

The two probability profiles in \eqref{eq:certificate-profiles} are the
normalized positive and negative parts of
\eqref{eq:certificate-annihilator}.  Any estimator receiving their common
four-jet data must therefore make error at least half of the $3/4$ window
gap on one of them.  This realizes the radius in
\eqref{eq:sharp-five-node-radius} without floating-point computation.
